% =====================================================================
%  vision.tex  --  Documento de Vision de LEIP
%  ARCHIVO UNICO: no necesita ningun otro archivo.
%  Compilar en Overleaf con pdfLaTeX.
% =====================================================================
\documentclass[11pt,a4paper]{article}

% ---------------------------------------------------------------------
%  DATOS DEL PROYECTO  (tocá solo esto para actualizar el documento)
% ---------------------------------------------------------------------
\newcommand{\proyectoNombre}{LEIP}
\newcommand{\proyectoNombreLargo}{LatAm Economic Intelligence Platform}
\newcommand{\proyectoUniversidad}{Universidad Cat\'olica del Uruguay}
\newcommand{\proyectoFacultad}{Facultad de Ingenier\'ia y Tecnolog\'ias}
\newcommand{\proyectoCarrera}{Ingenier\'ia en Inform\'atica}
\newcommand{\docVersion}{1.3}
\newcommand{\docTitulo}{Documento de Visi\'on}
\newcommand{\docFecha}{Septiembre 2026}
\newcommand{\equipoLista}{%
  Sebastian Forische \and
  Armando Hernandez de la Vega \and
  Luana Acu\~na \and
  Bruno Sebastian Olivera Viera Gomez%
}
\newcommand{\modMacro}{Macro Expectations Intelligence}
\newcommand{\modPrecios}{Prices \& Inflation Intelligence}
\newcommand{\modActividad}{Real Activity Intelligence}

% ---------------------------------------------------------------------
%  PREAMBULO  (estilos y comandos; normalmente no hace falta tocarlo)
% ---------------------------------------------------------------------
% =====================================================================
%  preamble.tex — Paquetes, estilo y comandos. Compartido por TODOS los docs.
% =====================================================================

\usepackage[utf8]{inputenc}
\usepackage[T1]{fontenc}
% microtype requiere fuentes escalables; va junto con lmodern.
\IfFileExists{lmodern.sty}{\usepackage{lmodern}\usepackage{microtype}}{}

% --- Idioma -------------------------------------------------------------
% Si babel-spanish está instalado (Overleaf, TeX Live completo) lo usa.
% Si no, cae a inglés y renombra los títulos a mano para no romper el build.
\IfFileExists{spanish.ldf}{%
  \usepackage[spanish,es-tabla,es-noquoting]{babel}%
}{%
  \usepackage[english]{babel}%
  \renewcommand{\contentsname}{\'Indice}%
  \renewcommand{\listfigurename}{\'Indice de figuras}%
  \renewcommand{\listtablename}{\'Indice de tablas}%
  \renewcommand{\tablename}{Tabla}%
  \renewcommand{\figurename}{Figura}%
  \renewcommand{\refname}{Referencias}%
  \renewcommand{\abstractname}{Resumen}%
}

% --- Página -------------------------------------------------------------
\usepackage[a4paper,margin=2.5cm,headheight=15pt]{geometry}
\usepackage{fancyhdr}
\usepackage{parskip}

% --- Tablas y listas ----------------------------------------------------
\usepackage{booktabs}
\usepackage{tabularx}
\usepackage{longtable}
\usepackage{array}
\usepackage{enumitem}
\setlist{noitemsep,topsep=4pt}

% --- Gráficos y color ---------------------------------------------------
\usepackage{graphicx}
\usepackage{xcolor}
\usepackage{tcolorbox}
\tcbuselibrary{breakable}

% --- Enlaces ------------------------------------------------------------
\usepackage[hidelinks]{hyperref}
\usepackage{url}
\urlstyle{same}

% --- Paleta -------------------------------------------------------------
\definecolor{leipAzul}{HTML}{1F4E79}
\definecolor{leipGris}{HTML}{5A5A5A}
\definecolor{leipAmbar}{HTML}{B8860B}
\definecolor{leipAmbarBg}{HTML}{FFF8E1}

% --- Encabezado y pie ---------------------------------------------------
\pagestyle{fancy}
\fancyhf{}
\fancyhead[L]{\small\color{leipGris}\proyectoNombre{} --- \docTitulo}
\fancyhead[R]{\small\color{leipGris} v\docVersion}
\fancyfoot[C]{\small\thepage}
\renewcommand{\headrulewidth}{0.4pt}

% --- Títulos de sección -------------------------------------------------
\usepackage{titlesec}
\titleformat{\section}{\Large\bfseries\color{leipAzul}}{\thesection}{0.6em}{}
\titleformat{\subsection}{\large\bfseries\color{leipAzul}}{\thesubsection}{0.6em}{}
\titleformat{\subsubsection}{\normalsize\bfseries\color{leipGris}}{\thesubsubsection}{0.6em}{}

% =====================================================================
%  COMANDOS PROPIOS
% =====================================================================

% --- Portada estándar ---------------------------------------------------
% Uso: \portada{\docTitulo}
\newcommand{\logoUniversidad}{img/logo-ucu.png}
\newcommand{\portada}[1]{%
  \begin{titlepage}
  \centering
  \vspace*{1.5cm}
  \IfFileExists{\logoUniversidad}{%
    \includegraphics[width=0.40\textwidth]{\logoUniversidad}\par
    \vspace{1.2cm}%
  }{%
    \fbox{\parbox[c][3cm][c]{0.5\textwidth}{\centering\color{leipAmbar}%
      \textbf{[Logo pendiente]}\\[4pt]\small\texttt{\logoUniversidad}}}\par
    \vspace{1.2cm}%
  }
  {\Large \proyectoUniversidad\par}
  {\large \proyectoFacultad\par}
  {\large \proyectoCarrera\par}
  \vspace{2.5cm}
  {\Huge\bfseries\color{leipAzul} #1\par}
  \vspace{0.8cm}
  {\Large \proyectoNombre{} --- \proyectoNombreLargo\par}
  \vspace{2cm}
  {\large Versi\'on \docVersion\par}
  {\large \docFecha\par}
  \vfill
  \begin{tabular}{c}
    Sebastian Forische \\
    Armando Hernandez de la Vega \\
    Luana Acu\~na \\
    Bruno Sebastian Olivera Viera Gomez \\
  \end{tabular}
  \vspace{1.5cm}
  \end{titlepage}
}

% --- Historial de revisiones -------------------------------------------
% Uso: \begin{historial} 1.0 & 26/06/2026 & Versión inicial & Armando \\ \end{historial}
% Nota: usa tabular con p{} y NO tabularx. tabularx necesita leer el cuerpo
% completo de la tabla de una sola vez, y no tolera que \begin y \end queden
% en macros distintas.
\newenvironment{historial}{%
  \section*{Historial de revisiones}
  \addcontentsline{toc}{section}{Historial de revisiones}
  \begin{tabular}{@{}l l p{0.52\textwidth} l@{}}
  \toprule
  \textbf{Versi\'on} & \textbf{Fecha} & \textbf{Descripci\'on} & \textbf{Autor} \\
  \midrule
}{%
  \bottomrule
  \end{tabular}
  \clearpage
}

% --- Marcador de contenido pendiente -----------------------------------
% Uso: \pendiente{resumen corto}{texto completo, puede traer listas}
%   - el resumen corto va al listado del final: debe ser texto plano
%   - el texto completo va al recuadro en el cuerpo del documento
% Para la entrega final: \pendientesfalse oculta los recuadros del PDF,
% pero el listado del final sigue mostrando qué quedó abierto.
\newif\ifpendientes
\pendientestrue

\makeatletter
\newcommand{\listofpendientes}{%
  \section*{Pendientes de este documento}
  \addcontentsline{toc}{section}{Pendientes de este documento}
  \@starttoc{pnd}%
}
\newcommand{\pendiente}[2]{%
  \addcontentsline{pnd}{section}{\protect\numberline{}#1}%
  \ifpendientes
    \begin{tcolorbox}[breakable,colback=leipAmbarBg,colframe=leipAmbar,
                      boxrule=0.6pt,left=6pt,right=6pt,top=4pt,bottom=4pt]
      \textbf{\color{leipAmbar}PENDIENTE.}\ #2
    \end{tcolorbox}
  \fi
}
\makeatother

% --- Identificadores de requerimientos ---------------------------------
% Se usan en requerimientos.tex; van acá para que la numeración sea única.
\newcounter{rfcounter}
\newcommand{\rf}[2]{% \rf{etiqueta}{texto}
  \refstepcounter{rfcounter}%
  \label{rf:#1}%
  \textbf{RF-\arabic{rfcounter}}\quad #2\par\smallskip
}
\newcounter{rnfcounter}
\newcommand{\rnf}[2]{%
  \refstepcounter{rnfcounter}%
  \label{rnf:#1}%
  \textbf{RNF-\arabic{rnfcounter}}\quad #2\par\smallskip
}

% --- Figura con fallback si la imagen todavía no existe ----------------
% Uso: \figuraLEIP{img/arquitectura.png}{Arquitectura de LEIP}{fig:arq}{0.9}
\newcommand{\figuraLEIP}[4]{%
  \begin{figure}[htbp]
  \centering
  \IfFileExists{#1}{\includegraphics[width=#4\textwidth]{#1}}{%
    \fbox{\parbox[c][4cm][c]{0.8\textwidth}{\centering\color{leipAmbar}%
      \textbf{[Figura pendiente]}\\[4pt]\small\texttt{#1}}}%
  }
  \caption{#2}
  \label{#3}
  \end{figure}
}

\begin{document}


\portada{\docTitulo}

\tableofcontents
\clearpage

\begin{historial}
1.0 & 26/06/2026 & Contenido inicial derivado del anteproyecto & Equipo \\
1.1 & 28/08/2026 & Ajustes de validaci\'on, fuentes, monetizaci\'on y arquitectura de ingesta & Equipo \\
1.2 & 28/08/2026 & Consolidaci\'on de secciones y depuraci\'on de pendientes & Equipo \\
1.3 & 08/09/2026 & Reorganizaci\'on del documento en torno a los aspectos t\'ecnicos y funcionales del producto. Se incorporan acr\'onimos, introducciones de secci\'on y criterios de priorizaci\'on. El an\'alisis de mercado, la validaci\'on comercial y el modelo de negocio se trasladan al Documento de Negocio & Equipo \\
\end{historial}

% =====================================================================
\section{Introducci\'on}

\subsection{Prop\'osito}

Este documento define el producto \proyectoNombre{} (\proyectoNombreLargo):
qu\'e es, qu\'e problema resuelve, a qui\'enes est\'a dirigido y qu\'e
capacidades generales ofrece. Se concentra en los aspectos t\'ecnicos y
funcionales del producto: describe \emph{qu\'e} necesitan los usuarios y
\emph{por qu\'e} existen esas necesidades, y deja el detalle de \emph{c\'omo}
el sistema las satisface a la Especificaci\'on de Requerimientos, la
Especificaci\'on de Casos de Uso y las Notas de Arquitectura y Dise\~no.

Los aspectos de mercado, la validaci\'on comercial del problema, la definici\'on
del negocio y los planes de contrataci\'on se desarrollan en el Documento de
Negocio y no forman parte de este documento.

Este documento constituye la versi\'on vigente de la definici\'on del producto.
No se considera cerrado: se actualiza cada vez que una decisi\'on de producto,
una nueva fuente de datos o un cambio de capacidades modifiquen su contenido, y
cada modificaci\'on queda asentada en el historial de revisiones.

\subsection{Alcance del documento}

Esta secci\'on delimita qu\'e cubre y qu\'e no cubre el documento, para que el
lector sepa d\'onde buscar la informaci\'on que no encuentre aqu\'i.

El documento cubre la definici\'on del producto y su perspectiva dentro del
entorno de trabajo de los usuarios, el problema que resuelve, los usuarios,
interesados y actores que interact\'uan con \'el, sus capacidades generales, las
asunciones y dependencias sobre las que se apoya, las restricciones t\'ecnicas,
operativas, legales y de integraci\'on que lo condicionan, y los criterios con
los que se prioriza su construcci\'on.

No cubre la especificaci\'on numerada de requerimientos, la descripci\'on
detallada de los casos de uso ni el dise\~no y la arquitectura de la soluci\'on,
que se desarrollan en sus documentos correspondientes. Tampoco cubre el
an\'alisis de mercado ni el modelo de negocio.

\subsection{Acr\'onimos}

A continuaci\'on se presentan los acr\'onimos utilizados a lo largo del
documento y una breve descripci\'on de cada uno. El Glosario contiene el
listado completo de t\'erminos y siglas del producto.

\begin{table}[htbp]
\centering
\caption{Acr\'onimos utilizados en el documento}
\label{tab:acronimos}
\begin{tabularx}{\textwidth}{l X}
\toprule
\textbf{Acr\'onimo} & \textbf{Descripci\'on} \\
\midrule
API     & \emph{Application Programming Interface}. Interfaz mediante la cual un sistema consulta o intercambia informaci\'on con otro de forma program\'atica. \\
Banxico & Banco de M\'exico. Banco central de M\'exico y fuente oficial de las expectativas de ese pa\'is. \\
BCB     & Banco Central do Brasil. Banco central de Brasil y fuente oficial del Bolet\'in Focus. \\
BCCh    & Banco Central de Chile. Banco central de Chile, fuente considerada para una etapa posterior. \\
BCU     & Banco Central del Uruguay. Banco central de Uruguay y fuente oficial de sus encuestas de expectativas. \\
BI      & \emph{Business Intelligence}. Conjunto de herramientas de an\'alisis y visualizaci\'on que los usuarios emplean en su entorno de trabajo. \\
CSV     & \emph{Comma-Separated Values}. Formato de archivo de texto plano en el que algunos organismos publican sus series. \\
HTML    & \emph{HyperText Markup Language}. Lenguaje de las p\'aginas web; algunas fuentes publican sus datos \'unicamente dentro de una p\'agina. \\
INE     & Instituto Nacional de Estad\'istica de Uruguay. Organismo emisor de indicadores de precios y actividad. \\
INEGI   & Instituto Nacional de Estad\'istica y Geograf\'ia de M\'exico. Organismo emisor de indicadores de precios y actividad. \\
JSON    & \emph{JavaScript Object Notation}. Formato estructurado de intercambio de datos, habitual en las respuestas de las API oficiales. \\
LLM     & \emph{Large Language Model}. Modelo de lenguaje utilizado como apoyo para interpretar fuentes heterog\'eneas y generar res\'umenes descriptivos a partir de m\'etricas ya calculadas. \\
MEF     & Ministerio de Econom\'ia y Finanzas de Uruguay. Organismo emisor de informaci\'on fiscal. \\
PDF     & \emph{Portable Document Format}. Formato de documento en el que varias fuentes publican sus informes y tablas. \\
SaaS    & \emph{Software as a Service}. Modelo bajo el cual el producto se ofrece como servicio accesible por red, sin instalaci\'on por parte del usuario. \\
SIE     & Sistema de Informaci\'on Econ\'omica de Banxico. Sistema desde el que se accede program\'aticamente a las series mexicanas. \\
URL     & \emph{Uniform Resource Locator}. Direcci\'on que identifica el recurso publicado por una fuente y se conserva como parte de la trazabilidad. \\
\bottomrule
\end{tabularx}
\end{table}

\subsection{Referencias}

A continuaci\'on se presentan las referencias utilizadas para la elaboraci\'on
del documento y el prop\'osito de cada una.

\begin{table}[htbp]
\centering
\caption{Referencias y prop\'osito de cada una}
\label{tab:referencias}
\begin{tabularx}{\textwidth}{X X}
\toprule
\textbf{Referencia} & \textbf{Para qu\'e se utiliza} \\
\midrule
Bolet\'in Focus --- Banco Central do Brasil. Expectativas de mercado, frecuencia semanal. &
Determina la cobertura y la frecuencia de actualizaci\'on del m\'odulo de expectativas para Brasil, y sirve de caso de referencia para el modelo com\'un de variable, unidad y horizonte. \\
\addlinespace
Encuesta de expectativas y Sistema de Informaci\'on Econ\'omica (SIE) --- Banco de M\'exico. Frecuencia quincenal o mensual. &
Determina la cobertura de M\'exico y aporta el caso de acceso program\'atico a las series, que condiciona el dise\~no de la ingesta. \\
\addlinespace
Encuestas de expectativas econ\'omicas --- Banco Central del Uruguay. Frecuencia mensual. &
Determina la cobertura de Uruguay y aporta el caso de publicaci\'on en archivos descargables, con la variabilidad de formato que la normalizaci\'on debe absorber. \\
\addlinespace
Encuesta de Expectativas Econ\'omicas --- Banco Central de Chile. Frecuencia mensual. &
Se utiliza para delimitar el alcance geogr\'afico inicial y para ilustrar por qu\'e las condiciones de uso se verifican fuente por fuente antes de incorporar un pa\'is. \\
\addlinespace
Encuesta de relevamiento del problema.
\url{https://docs.google.com/forms/d/e/1FAIpQLSe4jNRspwc_cyIJii9npmMjpXSXsAxOBzT6vkGCEi5vW80DEQ/viewform} &
Instrumento con el que se releva la frecuencia y el costo operativo de las tareas que el producto automatiza. Sus resultados sustentan la definici\'on del problema y la priorizaci\'on de las capacidades. \\
\bottomrule
\end{tabularx}
\end{table}

Las referencias anteriores corresponden al m\'odulo inicial de expectativas.
Durante la incorporaci\'on de cada conjunto de datos se mantiene un registro de
fuentes que incluye organismo, URL, fecha de consulta, m\'etodo de acceso,
condiciones de uso y requisitos de atribuci\'on. La evaluaci\'on se realiza
fuente por fuente y no mediante una licencia \'unica por pa\'is.

% =====================================================================
\section{Definici\'on del problema}

Esta secci\'on describe la situaci\'on que da origen al producto: qu\'e
dificultad enfrentan hoy quienes trabajan con informaci\'on econ\'omica oficial
de Am\'erica Latina, a qui\'enes afecta, qu\'e consecuencias tiene y qu\'e
deber\'ia resolver una soluci\'on adecuada.

\subsection{Problema identificado}

La informaci\'on econ\'omica oficial de Am\'erica Latina es p\'ublica, pero no
es directamente comparable. Cada pa\'is la publica a trav\'es de organismos
distintos ---bancos centrales, institutos de estad\'istica y ministerios---, en
formatos heterog\'eneos ---API, Excel, CSV, PDF, HTML o JSON--- y bajo nombres,
unidades, frecuencias, coberturas y metodolog\'ias que no siempre son
equivalentes entre s\'i.

Como consecuencia, quien necesita analizar la regi\'on repite un mismo trabajo
previo antes de poder analizar: ingresar a cada organismo, localizar la
publicaci\'on vigente, descargar el archivo, interpretar su estructura,
homogeneizar variables y unidades, y reconciliar las series entre pa\'ises. Ese
trabajo se rehace cada vez que una fuente publica una actualizaci\'on, y su
resultado queda en planillas o \emph{scripts} propios que alguien debe
mantener.

A esto se suma que las fuentes revisan sus datos. Una expectativa o un
indicador publicado en un per\'iodo puede corregirse en publicaciones
posteriores, y la versi\'on anterior deja de estar disponible en el sitio del
organismo. Reconstruir esa historia exige haber conservado cada versi\'on en el
momento en que fue publicada, algo que los procesos manuales rara vez sostienen
de forma sistem\'atica.

El resultado es tiempo operativo recurrente dedicado a tareas que no son
anal\'iticas, riesgo de error al comparar pa\'ises cuyas series no son
equivalentes, e imposibilidad pr\'actica de seguir c\'omo evolucionaron las
mediciones a lo largo del tiempo.

\subsection{Enunciado formal}

La tabla \ref{tab:problema} resume el problema en sus cuatro dimensiones: en
qu\'e consiste, a qui\'enes alcanza, qu\'e consecuencias produce y qu\'e
deber\'ia lograr una soluci\'on adecuada.

\begin{table}[htbp]
\centering
\begin{tabularx}{\textwidth}{>{\bfseries}l X}
\toprule
El problema de & la fragmentaci\'on de la informaci\'on econ\'omica oficial de
Am\'erica Latina entre m\'ultiples organismos, formatos, unidades, frecuencias
y metodolog\'ias no equivalentes. \\
\addlinespace
Afecta a & analistas macroecon\'omicos, equipos de investigaci\'on,
tesorer\'ias, consultoras, fondos de inversi\'on, brokers, acad\'emicos e
investigadores que trabajan de forma recurrente con datos de la regi\'on. \\
\addlinespace
Cuyo impacto es & tiempo operativo recurrente dedicado a buscar, descargar,
limpiar y reconciliar datos; riesgo de error en la comparaci\'on entre
pa\'ises; imposibilidad pr\'actica de seguir las revisiones hist\'oricas de
cada serie; y dependencia de planillas o \emph{scripts} internos que hay que
mantener. \\
\addlinespace
Una soluci\'on exitosa & entregar\'ia informaci\'on oficial ya normalizada y
comparable entre pa\'ises, con trazabilidad de origen y versi\'on, seguimiento
expl\'icito de revisiones, y m\'etricas y cruces listos para el an\'alisis, sin
que el usuario deba mantener procesos propios de descarga y limpieza. \\
\bottomrule
\end{tabularx}
\caption{Enunciado formal del problema}
\label{tab:problema}
\end{table}

% =====================================================================
\section{Oportunidad y contexto}

Esta secci\'on explica de d\'onde surge el producto, c\'omo organiza el dominio
econ\'omico que cubre y sobre qu\'e fuentes oficiales trabaja en una primera
etapa.

\subsection{Origen de la idea}

La idea de \proyectoNombre{} surge de una necesidad observada en el trabajo con
informaci\'on econ\'omica y financiera de Am\'erica Latina. Los organismos
oficiales publican datos relevantes, pero \'estos se encuentran dispersos entre
m\'ultiples fuentes, formatos, unidades, nombres, frecuencias y metodolog\'ias.

En la pr\'actica, esa fragmentaci\'on obliga a analistas, investigadores y
equipos financieros a invertir tiempo en buscar publicaciones, descargar
archivos, limpiar datos, mantener planillas o \emph{scripts} internos y adaptar
las series antes de poder analizarlas y compararlas entre pa\'ises.

A partir de esta observaci\'on surge \proyectoNombre{}, una plataforma SaaS de
inteligencia econ\'omica que transforma informaci\'on oficial, dispersa y
heterog\'enea en m\'etricas comparables, trazables y listas para el an\'alisis.
La infraestructura de ingesta y normalizaci\'on funciona como base habilitante,
mientras que la capa de inteligencia, la navegaci\'on comparable y la
distribuci\'on recurrente de se\~nales son las que resuelven directamente la
necesidad del usuario.

\subsection{Organizaci\'on del dominio}

El producto organiza la informaci\'on en m\'odulos de expectativas
macroecon\'omicas, precios e inflaci\'on, actividad econ\'omica, riesgo fiscal,
sector externo y condiciones monetarias y financieras, dejando los activos de
mercado para una etapa posterior. Los m\'odulos pueden relacionarse entre s\'i
para comparar, por ejemplo, inflaci\'on esperada frente a observada, riesgo
fiscal frente a tasas soberanas, o reservas frente a presi\'on cambiaria.

\subsection{Cobertura inicial y fuentes}

Para una primera etapa, \proyectoNombre{} trabaja sobre informaci\'on oficial de
Brasil, M\'exico y Uruguay, y desarrolla los m\'odulos de expectativas
macroecon\'omicas, precios e inflaci\'on y actividad econ\'omica. Chile queda
fuera de esa primera etapa: su incorporaci\'on depende de obtener autorizaci\'on
expresa para el uso previsto de sus datos.

La tabla \ref{tab:fuentes} resume, para cada pa\'is, la fuente de expectativas
considerada, su frecuencia de publicaci\'on y su situaci\'on respecto del
alcance inicial.

\begin{table}[htbp]
\centering
\caption{Fuentes iniciales del m\'odulo de expectativas}
\label{tab:fuentes}
\begin{tabularx}{\textwidth}{l l X}
\toprule
\textbf{Pa\'is} & \textbf{Fuente considerada} & \textbf{Frecuencia y situaci\'on} \\
\midrule
Brasil  & Bolet\'in Focus (BCB)            & Semanal. Incluida en la primera etapa. \\
M\'exico & Encuesta de expectativas / SIE (Banxico) & Quincenal o mensual. Incluida en la primera etapa. \\
Uruguay & Encuestas de expectativas (BCU)  & Mensual. Incluida en la primera etapa. \\
Chile   & Encuesta de Expectativas Econ\'omicas (BCCh) & Mensual. Fuera de la primera etapa; su incorporaci\'on depende de la autorizaci\'on de uso del organismo. \\
\bottomrule
\end{tabularx}
\end{table}

La tabla \ref{tab:fuentes} presenta solamente las fuentes iniciales del m\'odulo
de expectativas. Los m\'odulos de precios, actividad y los cruces anal\'iticos
incorporan tambi\'en organismos estad\'isticos, ministerios y otras fuentes
oficiales, como INE, INEGI, MEF y la Secretar\'ia de Econom\'ia. La procedencia
y las condiciones de uso se registran para cada conjunto de datos; ninguna
instituci\'on representa por s\'i sola todas las fuentes de un pa\'is.

% =====================================================================
\section{Definici\'on del producto}

Esta secci\'on presenta qu\'e es \proyectoNombre{}, qu\'e resuelve, c\'omo se
inserta en el entorno de trabajo de sus usuarios y c\'omo garantiza que cada
dato que entrega pueda rastrearse hasta su origen.

\subsection{Definici\'on}

\proyectoNombre{} es una plataforma SaaS que centraliza informaci\'on
econ\'omica oficial de Am\'erica Latina y la entrega normalizada, comparable
entre pa\'ises y trazable hasta la publicaci\'on que le dio origen.

El producto resuelve el trabajo previo descrito en la secci\'on anterior. En
lugar de que cada usuario recorra los organismos, descargue archivos y
homogeneice series, la plataforma ejecuta esa ingesta de forma automatizada y
continua, conserva cada versi\'on publicada y expone el resultado bajo un modelo
com\'un de pa\'is, variable, unidad, frecuencia y horizonte. El usuario deja de
mantener procesos propios de descarga, limpieza y actualizaci\'on.

Sobre esa base, el producto ofrece tres capacidades generales:

\begin{itemize}
  \item \textbf{Consulta y comparaci\'on.} Acceso a series de varios pa\'ises
        bajo criterios homog\'eneos, con hist\'oricos completos y con las
        revisiones sucesivas de cada medici\'on conservadas y consultables.
  \item \textbf{Inteligencia sobre los datos.} M\'etricas determin\'isticas
        ---cambio, momentum, divergencia regional, percentiles hist\'oricos,
        sorpresa, persistencia y distancia frente a metas--- y cruces entre
        m\'odulos que relacionan lo esperado con lo observado.
  \item \textbf{Distribuci\'on e integraci\'on.} Exportaciones auditables,
        alertas, \emph{briefings} verificables y acceso program\'atico para que
        las organizaciones incorporen los datos y las m\'etricas a sus propios
        procesos.
\end{itemize}

Toda salida del producto ---un dato, una m\'etrica, un gr\'afico o una
exportaci\'on--- conserva la referencia al organismo, el conjunto de datos, la
fecha de publicaci\'on y la versi\'on de los que proviene.

\subsection{Perspectiva del producto}

\proyectoNombre{} es un sistema autocontenido que se ubica entre las fuentes
oficiales y las herramientas con las que los usuarios ya trabajan. No reemplaza
esas herramientas ni el criterio anal\'itico de quien las usa: se inserta antes
del an\'alisis y entrega el insumo ya preparado, de modo que el usuario pueda
comenzar por la interpretaci\'on y no por la recolecci\'on.

Su funci\'on es operar de forma continua un ciclo que hoy los usuarios ejecutan
a mano. La plataforma identifica la fuente, extrae la informaci\'on desde
formatos diversos, valida su contenido, normaliza variables, unidades, fechas,
frecuencias y metodolog\'ias, y almacena el resultado en un Data Hub organizado
por pa\'is, variable, per\'iodo, unidad, fuente, fecha de publicaci\'on y
versi\'on. Sobre ese repositorio se calculan las m\'etricas y los cruces que
constituyen la capa de inteligencia.

El producto se relaciona con tres tipos de sistemas externos:

\begin{itemize}
  \item \textbf{Fuentes oficiales.} Portales, API y archivos publicados por
        bancos centrales, institutos de estad\'istica y ministerios de los
        pa\'ises cubiertos. Son la \'unica entrada de datos econ\'omicos y
        determinan la frecuencia de actualizaci\'on de cada serie.
  \item \textbf{Herramientas de los usuarios.} Navegador web, planillas de
        c\'alculo, entornos de an\'alisis como R o Python y herramientas de BI.
        El producto llega a ellas mediante la aplicaci\'on web, las
        exportaciones auditables y el acceso program\'atico.
  \item \textbf{Servicios de terceros.} Un proveedor externo de cobros, que
        expone \emph{checkout} alojado y notificaciones por \emph{webhook}, y
        los mecanismos de autenticaci\'on y control de acceso de las
        organizaciones cliente.
\end{itemize}

Desde el punto de vista funcional, el valor del producto est\'a en entregar
comparabilidad regional lista para usar: el usuario obtiene en una consulta lo
que antes requer\'ia recorrer varios organismos, y puede sostener un seguimiento
peri\'odico sin rehacer el trabajo en cada publicaci\'on. Desde el punto de
vista t\'ecnico, el valor est\'a en el versionado y la trazabilidad: la
plataforma conserva la historia de cada medici\'on y el origen de cada n\'umero,
lo que permite auditar un resultado y reproducirlo, algo que los procesos
manuales no garantizan.

Como primera hip\'otesis t\'ecnica, se definir\'a un boceto de arquitectura de
agentes para los procesos de extracci\'on, cuyo objetivo ser\'a coordinar la
obtenci\'on, interpretaci\'on y validaci\'on de fuentes heterog\'eneas
manteniendo controles determin\'isticos de calidad y trazabilidad. La
composici\'on de los agentes, sus herramientas y su orquestaci\'on se
especifican en las Notas de Arquitectura y Dise\~no.

\figuraLEIP{img/arquitectura-contexto.png}{Vista de contexto de \proyectoNombre{}}{fig:contexto}{0.9}

La figura \ref{fig:contexto} presenta a \proyectoNombre{} como una caja \'unica
y muestra las fuentes oficiales, el proveedor de pagos y los sistemas
consumidores como entidades externas. El diagrama t\'ecnico completo, incluidas
las capas internas de procesamiento, se desarrolla en las Notas de
Arquitectura.

\subsection{Trazabilidad metodol\'ogica}

Un aspecto central del producto es la trazabilidad metodol\'ogica. Cada dato,
m\'etrica y exportaci\'on conserva su organismo emisor, conjunto de datos, URL,
fecha de publicaci\'on, versi\'on y criterio de procesamiento.

\proyectoNombre{} distingue expl\'icitamente entre el dato oficial y la
elaboraci\'on propia. Las series que los organismos publican se presentan
siempre con atribuci\'on y enlace a la fuente original. Las m\'etricas,
comparaciones y lecturas anal\'iticas calculadas por la plataforma se
identifican como \emph{Elaboraci\'on LEIP con datos de} la fuente
correspondiente, y no se presentan en ning\'un caso como indicadores oficiales
del organismo emisor.

Las condiciones de uso se eval\'uan por organismo y conjunto de datos, no por
pa\'is ni exclusivamente por banco central. Esta trazabilidad es a la vez una
capacidad funcional del producto y la principal mitigaci\'on de las
restricciones descritas en la secci\'on \ref{sec:restricciones}.

% =====================================================================
\section{Usuarios e interesados}

Esta secci\'on identifica a quienes interact\'uan con el producto o se ven
afectados por \'el: los interesados que condicionan su funcionamiento, los
perfiles de usuario que lo utilizan, el entorno en el que trabajan y las
necesidades que el producto debe atender.

\subsection{Interesados}

La tabla \ref{tab:interesados} presenta los interesados que participan en el
funcionamiento del producto y el criterio con el que cada uno juzga su
\'exito.

\begin{table}[htbp]
\centering
\caption{Interesados principales}
\label{tab:interesados}
\begin{tabularx}{\textwidth}{l X X}
\toprule
\textbf{Interesado} & \textbf{Participaci\'on} & \textbf{Criterio de \'exito} \\
\midrule
Organismos emisores & Publican datos, metodolog\'ias y condiciones de uso & Atribuci\'on correcta y uso conforme a sus condiciones \\
\addlinespace
Usuarios de la plataforma & Consultan, comparan y exportan informaci\'on en su trabajo diario & Reducci\'on del trabajo manual y confianza en el dato entregado \\
\addlinespace
Organizaciones cliente & Incorporan datos y m\'etricas a sus procesos internos & Informaci\'on confiable, integrable y auditable \\
\addlinespace
Administradores de la organizaci\'on & Gestionan usuarios, roles y permisos de su instituci\'on & Control efectivo del acceso a m\'odulos, pa\'ises y funcionalidades \\
\addlinespace
Responsables t\'ecnicos de integraci\'on & Consumen la plataforma desde sistemas propios & Mecanismos de integraci\'on estables, documentados y versionados \\
\bottomrule
\end{tabularx}
\end{table}

\subsection{Perfiles de usuario}

Los perfiles de la tabla \ref{tab:perfiles} son la base de la que se derivan los
actores y los casos de uso del sistema. Para cada uno se indica su objetivo
principal y las capacidades del producto que utiliza con mayor frecuencia.

\begin{table}[htbp]
\centering
\caption{Perfiles de usuario}
\label{tab:perfiles}
\begin{tabularx}{\textwidth}{l X X}
\toprule
\textbf{Perfil} & \textbf{Objetivo principal} & \textbf{Uso esperado} \\
\midrule
Analista macroecon\'omico & Seguir variables y expectativas regionales & Comparadores, hist\'oricos, revisiones y m\'etricas \\
\addlinespace
Equipo de investigaci\'on & Producir an\'alisis recurrente y verificable & Radar, consultas guardadas, exportaciones e integraci\'on program\'atica \\
\addlinespace
Tesorer\'ia o fondo & Incorporar contexto macro a decisiones internas & Alertas, cruces, escenarios y series comparables \\
\addlinespace
Consultora o broker & Preparar informes para clientes & Gr\'aficos, \emph{briefings} y exportaciones auditables \\
\addlinespace
Acad\'emico o investigador & Analizar series regionales con trazabilidad & Hist\'oricos, metadatos y descarga estructurada \\
\addlinespace
Administrador institucional & Gestionar el acceso de su organizaci\'on & Usuarios, roles, permisos y \emph{entitlements} \\
\bottomrule
\end{tabularx}
\end{table}

\subsection{Entorno de los usuarios}

Los usuarios de \proyectoNombre{} trabajan en \'areas de an\'alisis econ\'omico
y financiero: mesas de tesorer\'ia, equipos de \emph{research}, consultoras,
\'areas de estudios y grupos de investigaci\'on. Su entorno habitual combina el
navegador web con planillas de c\'alculo, entornos de an\'alisis como R o Python
y herramientas de BI como Power BI.

En ese entorno realizan tareas recurrentes cuyo ritmo lo marca el calendario de
publicaci\'on de los organismos: revisar qu\'e public\'o cada fuente, actualizar
las series que mantienen, comparar la evoluci\'on entre pa\'ises, contrastar lo
esperado con lo observado y preparar informes o seguimientos peri\'odicos, sea
para clientes externos o para decisiones internas.

De esas tareas se desprenden tres necesidades. La primera es disponer del dato
ya normalizado y comparable, sin ejecutar cada vez el proceso de b\'usqueda,
descarga y limpieza. La segunda es poder justificar el origen de cada n\'umero
que publican, incluso cuando la fuente revis\'o su propia serie despu\'es de la
consulta. La tercera es llevar esa informaci\'on a sus propias herramientas sin
pasos manuales intermedios.

El producto atiende esas necesidades en dos modos de uso. Quien trabaja de
forma individual accede a trav\'es de la aplicaci\'on web, consulta y compara
en el \emph{dashboard} y se lleva lo que necesita mediante exportaciones
auditables. Las organizaciones con procesos propios consumen los mismos datos y
m\'etricas de forma program\'atica e integran el resultado en sus flujos
internos sin intervenci\'on manual.

\subsection{Necesidades clave}

La tabla \ref{tab:necesidades} relaciona cada necesidad identificada en los
usuarios con la capacidad del producto que la atiende.

\begin{table}[htbp]
\centering
\caption{Necesidades de los usuarios y c\'omo las atiende \proyectoNombre{}}
\label{tab:necesidades}
\begin{tabularx}{\textwidth}{X X}
\toprule
\textbf{Necesidad} & \textbf{C\'omo la atiende \proyectoNombre{}} \\
\midrule
Dejar de invertir tiempo en tareas repetitivas de descarga, limpieza y
normalizaci\'on &
Ingesta automatizada y normalizaci\'on centralizada de las fuentes oficiales \\
\addlinespace
Comparar pa\'ises, variables y horizontes bajo criterios homog\'eneos &
Modelo com\'un de pa\'is, variable, unidad, frecuencia y horizonte \\
\addlinespace
Saber qu\'e cambi\'o respecto de mediciones anteriores, no s\'olo el dato puntual &
Conservaci\'on de revisiones y m\'etricas de cambio sobre el hist\'orico \\
\addlinespace
Poder justificar de d\'onde sali\'o cada n\'umero &
Trazabilidad de fuente, fecha, versi\'on y criterio metodol\'ogico en cada dato,
m\'etrica y exportaci\'on \\
\addlinespace
Producir reportes y seguimientos recurrentes sin rehacer el trabajo cada vez &
Exportaciones auditables, consultas guardadas y radar peri\'odico \\
\addlinespace
Detectar qu\'e movimientos merecen atenci\'on &
M\'etricas de cambio, sorpresa, divergencia regional y distancia frente a metas \\
\addlinespace
Incorporar la informaci\'on a herramientas y procesos propios &
Exportaciones en formatos estructurados y acceso program\'atico autenticado \\
\bottomrule
\end{tabularx}
\end{table}

% =====================================================================
\section{Capacidades del producto}

Las capacidades siguientes constituyen el alcance funcional previsto para una
primera etapa del producto:

\begin{itemize}
  \item \textbf{Pa\'ises cubiertos:} Brasil, M\'exico y Uruguay, utilizando
        fuentes oficiales sujetas a las condiciones de uso de cada organismo.
  \item \textbf{M\'odulos:} \modMacro{}, \modPrecios{} y \modActividad{}.
  \item \textbf{M\'etricas determin\'isticas:} cambios, momentum, divergencias
        regionales, percentiles hist\'oricos, distancia frente a metas,
        sorpresas y persistencia.
  \item \textbf{Cruces:} inflaci\'on esperada frente a inflaci\'on observada, y
        crecimiento esperado frente a actividad efectiva.
  \item \textbf{Aplicaci\'on web:} \emph{dashboards}, comparadores, fichas por
        pa\'is, radar, alertas, \emph{briefings} verificables y exportaciones
        auditables.
  \item \textbf{Control de acceso:} autenticaci\'on, roles, permisos y
        \emph{entitlements} para controlar el acceso a m\'odulos, pa\'ises y
        funcionalidades.
  \item \textbf{Integraci\'on program\'atica:} acceso autenticado para
        organizaciones que necesiten incorporar datos y m\'etricas a sus propios
        procesos.
\end{itemize}

% =====================================================================
\section{Asunciones y dependencias}

El producto se define bajo las asunciones que se enumeran a continuaci\'on. Si
alguna de ellas no se cumple, el alcance previsto para la primera etapa debe
revisarse.

\begin{itemize}
  \item Los organismos oficiales mantienen la publicaci\'on de las series
        cubiertas, con frecuencia y cobertura equivalentes a las actuales.
  \item Los formatos de publicaci\'on se mantienen razonablemente estables, o
        sus cambios son detectables por validaci\'on autom\'atica.
  \item Cada organismo y conjunto de datos habilita el uso previsto por
        \proyectoNombre{} conforme a sus condiciones espec\'ificas. Esta
        verificaci\'on se mantiene en la matriz de fuentes y no se presume a
        partir del pa\'is de origen (ver secci\'on \ref{sec:restricciones}).
  \item Existe un proveedor externo de cobros que expone \emph{checkout} alojado
        y notificaciones por \emph{webhook}.
  \item La infraestructura del proveedor de nube est\'a disponible con la
        capacidad y las condiciones de servicio previstas.
\end{itemize}

Cada asunci\'on se vincula con su riesgo correspondiente en la Lista de
Riesgos, donde se registran probabilidad, impacto y mitigaci\'on.

% =====================================================================
\section{Restricciones del producto}
\label{sec:restricciones}

Esta secci\'on re\'une las condiciones t\'ecnicas, operativas, legales y de
integraci\'on que limitan lo que el producto puede hacer con los datos que
procesa y con la informaci\'on que entrega.

La restricci\'on principal es de origen legal y proviene de las condiciones de
uso aplicables a cada fuente oficial. Un segundo riesgo regulatorio, el de ser
interpretado como asesor\'ia financiera, es bajo por dise\~no:
\proyectoNombre{} entrega informaci\'on agregada y descriptiva, sin
recomendaciones personalizadas de inversi\'on.

No existe una licencia \'unica para los datos de un pa\'is ni una fuente
\'unica que lo represente. Brasil, M\'exico, Uruguay y las futuras
incorporaciones utilizan bancos centrales, institutos de estad\'istica,
ministerios y otros organismos. Por ello, cada conjunto de datos debe contar
con un registro de procedencia que incluya instituci\'on, URL, fecha de
consulta, m\'etodo de acceso, condiciones de uso y requisitos de atribuci\'on.

Seg\'un su origen, la informaci\'on que maneja el producto se somete a
restricciones distintas:

\begin{itemize}
  \item \textbf{Dato oficial publicado:} puede exponerse cuando las condiciones
        del organismo lo permitan, y conserva siempre la referencia a la fuente
        original.
  \item \textbf{Hist\'orico consolidado y versionado:} su almacenamiento y su
        presentaci\'on dependen de que las condiciones de cada conjunto de datos
        habiliten conservar y mostrar versiones anteriores.
  \item \textbf{M\'etricas y comparaciones propias:} son elaboraciones
        calculadas sobre datos trazables. Deben identificarse como
        \emph{Elaboraci\'on LEIP con datos de} y no presentarse como
        indicadores oficiales del organismo emisor.
\end{itemize}

Las restricciones se aplican fuente por fuente. Una limitaci\'on de un organismo
afecta solamente a los conjuntos de datos obtenidos de ese organismo y no
bloquea autom\'aticamente todas las fuentes del pa\'is. Del mismo modo, que una
serie sea p\'ublicamente accesible no sustituye la verificaci\'on de sus
condiciones de uso para hist\'oricos, exportaciones o redistribuci\'on.

Otras restricciones del producto:

\begin{itemize}
  \item \proyectoNombre{} no almacena ni procesa directamente datos de tarjetas
        de pago; el cobro se delega \'integramente a un proveedor externo.
  \item Las fuentes privadas que aporten las organizaciones cliente se procesan
        en entornos aislados y no se mezclan con la informaci\'on de otras
        organizaciones.
  \item La plataforma no genera datos propios ni completa vac\'ios de las
        fuentes: si un organismo no publica una serie, esa serie no est\'a
        disponible en el producto.
\end{itemize}

El cumplimiento de estas restricciones se apoya en la matriz de fuentes, un
registro que se mantiene actualizado durante toda la vida del producto. Para
cada conjunto de datos deja asentados la instituci\'on emisora, la URL, la
fecha de consulta, las condiciones de uso, la atribuci\'on requerida y qu\'e
est\'a permitido hacer con \'el en materia de almacenamiento, transformaci\'on,
hist\'oricos, exportaci\'on y redistribuci\'on. Ning\'un conjunto de datos se
incorpora al producto sin su entrada correspondiente en esa matriz.

% =====================================================================
\section{Rangos de calidad}

Los rangos de calidad son los umbrales medibles que debe alcanzar el producto
para considerarse aceptable en operaci\'on. Una vez definidos, se traducen en
requerimientos no funcionales numerados en la Especificaci\'on de
Requerimientos.

\pendiente{Definir los rangos de calidad y sus umbrales medibles}{Secci\'on a completar. Definir los umbrales medibles que hacen
aceptable al producto. Dimensiones candidatas:
\begin{itemize}
  \item \textbf{Frescura del dato:} m\'aximo de horas entre la publicaci\'on de
        la fuente y su disponibilidad en la plataforma.
  \item \textbf{Autonom\'ia de la ingesta:} porcentaje de ejecuciones que se
        completan sin intervenci\'on manual.
  \item \textbf{Cobertura de trazabilidad:} porcentaje de datos y m\'etricas con
        fuente, fecha y versi\'on completas (objetivo razonable: 100\%).
  \item \textbf{Tasa de cuarentena:} porcentaje aceptable de registros retenidos
        por ejecuci\'on.
  \item \textbf{Tiempo de respuesta} de una consulta comparativa en el
        \emph{dashboard}.
  \item \textbf{Disponibilidad} del servicio.
\end{itemize}}

% =====================================================================
\section{Criterios de priorizaci\'on}

El orden en que se construyen las capacidades del producto no es arbitrario:
responde a los criterios que se enuncian a continuaci\'on, aplicados en el orden
en que se presentan.

\begin{itemize}
  \item \textbf{Dependencia t\'ecnica.} Una capacidad no puede construirse antes
        que aquellas de las que depende. La ingesta y la normalizaci\'on
        preceden a toda la capa de inteligencia, porque sin datos comparables no
        hay m\'etrica posible.
  \item \textbf{Cobertura del n\'ucleo del problema.} Tiene prioridad lo que
        resuelve directamente la dificultad descrita en la definici\'on del
        problema, frente a lo que agrega alcance sobre un problema ya resuelto.
  \item \textbf{Habilitaci\'on de capacidades posteriores.} Se adelanta aquello
        que desbloquea varias funciones siguientes; un m\'odulo que habilita
        cruces vale m\'as, en orden de construcci\'on, que uno que se agota en
        s\'i mismo.
  \item \textbf{Sostenimiento del uso recurrente.} Entre dos capacidades
        equivalentes, se prioriza la que sostiene el uso peri\'odico del
        producto sobre la que se utiliza de forma puntual.
  \item \textbf{Relaci\'on entre esfuerzo y valor.} A igualdad de aporte, se
        elige la alternativa de menor esfuerzo de construcci\'on y
        mantenimiento.
\end{itemize}

Aplicando estos criterios, el orden de construcci\'on previsto es el siguiente:

\begin{enumerate}
  \item M\'odulo \modMacro{} sobre los tres pa\'ises cubiertos, como n\'ucleo
        del producto.
  \item M\'odulos \modPrecios{} y \modActividad{}, que habilitan los dos cruces
        iniciales.
  \item Radar, alertas y exportaciones, que sostienen el uso recurrente.
  \item Integraci\'on program\'atica y aislamiento multi-\emph{tenant},
        requeridos por las organizaciones que incorporan el producto a sus
        sistemas.
\end{enumerate}

Quedan fuera de la primera etapa y se abordan en etapas posteriores: los
m\'odulos de riesgo fiscal, sector externo y condiciones monetarias y
financieras; los activos de mercado; y la incorporaci\'on de Chile,
condicionada a obtener autorizaci\'on expresa del BCCh.

% =====================================================================
\section{Requerimientos de documentaci\'on}

El producto est\'a acompa\~nado por la documentaci\'on necesaria para que sus
datos, m\'etricas y mecanismos de acceso puedan utilizarse y auditarse. Los
documentos que se incluir\'an son los siguientes:

\begin{itemize}
  \item \textbf{Gu\'ia de usuario:} navegaci\'on, consultas, comparaciones,
        alertas y exportaciones.
  \item \textbf{Cat\'alogo de fuentes:} organismo, conjunto de datos, URL,
        frecuencia, fecha de consulta y condiciones de uso.
  \item \textbf{Diccionario de datos:} variables, unidades, frecuencias,
        horizontes, campos de trazabilidad y reglas de normalizaci\'on.
  \item \textbf{Notas metodol\'ogicas:} definici\'on, f\'ormula,
        interpretaci\'on, limitaciones y fuentes de cada m\'etrica propia.
  \item \textbf{Documento de integraci\'on con el sistema:} mecanismos de
        comunicaci\'on e integraci\'on que el producto pone a disposici\'on de
        sistemas externos, con sus condiciones de acceso, autenticaci\'on,
        par\'ametros, esquemas de respuesta, l\'imites de uso y ejemplos, con
        independencia de la tecnolog\'ia con la que se implementen.
  \item \textbf{Historial de cambios:} nuevas fuentes, revisiones
        metodol\'ogicas y modificaciones relevantes del producto.
\end{itemize}

Para la primera etapa son obligatorios la gu\'ia b\'asica, el cat\'alogo de
fuentes, el diccionario de datos y las notas metodol\'ogicas de las m\'etricas
implementadas. El documento de integraci\'on se publica cuando se habilite el
acceso desde sistemas externos. Todos ellos se mantienen actualizados a medida
que el producto evoluciona.

% =====================================================================
\clearpage
\listofpendientes

\end{document}
